\documentclass[11pt]{article}

\usepackage[margin=0.82in,headheight=22pt]{geometry}
\usepackage{amsmath,amssymb,mathtools,bm}
\usepackage{fontspec}
\setmainfont{TeX Gyre Pagella}
\setsansfont{TeX Gyre Heros}
\setmonofont{Courier New}[Scale=0.86]
\usepackage{microtype}
\usepackage{xcolor}
\usepackage{booktabs,longtable,array,colortbl}
\usepackage{enumitem}
\usepackage{fancyhdr}
\usepackage{titlesec}
\usepackage{hyperref}
\usepackage{bookmark}

\definecolor{navy}{RGB}{24,55,95}
\definecolor{teal}{RGB}{20,112,104}
\definecolor{softblue}{RGB}{239,245,251}
\definecolor{softgreen}{RGB}{239,248,245}
\definecolor{rulegray}{RGB}{205,214,222}

\hypersetup{
  colorlinks=true,
  linkcolor=navy,
  urlcolor=teal,
  pdftitle={STAT452 Take-home Mid-term Retake: Comprehensive Solutions},
  pdfauthor={STAT452}
}

\pagestyle{fancy}
\fancyhf{}
\lhead{\sffamily\small STAT452 Take-home Mid-term Retake}
\rhead{\sffamily\small Comprehensive Solutions}
\cfoot{\sffamily\thepage}
\renewcommand{\headrulewidth}{0.4pt}
\renewcommand{\headrule}{%
  \hbox to\headwidth{\color{navy!45}\leaders\hrule height \headrulewidth\hfill}}

\titleformat{\section}
  {\sffamily\LARGE\bfseries\color{navy}}{\thesection}{0.65em}{}
\titleformat{\subsection}
  {\sffamily\Large\bfseries\color{teal}}{\thesubsection}{0.65em}{}
\titleformat{\subsubsection}
  {\sffamily\large\bfseries\color{navy!85}}{\thesubsubsection}{0.65em}{}
\titlespacing*{\section}{0pt}{1.9em}{0.7em}
\titlespacing*{\subsection}{0pt}{1.35em}{0.45em}
\titlespacing*{\subsubsection}{0pt}{1.05em}{0.35em}

\setlength{\parindent}{0pt}
\setlength{\parskip}{0.48em}
\setlist{leftmargin=1.55em,itemsep=0.25em,topsep=0.32em}
\renewcommand{\arraystretch}{1.22}
\allowdisplaybreaks

\newcommand{\trans}{\mathsf T}
\newcommand{\E}{\operatorname{E}}
\newcommand{\Var}{\operatorname{Var}}
\newcommand{\Cov}{\operatorname{Cov}}
\newcommand{\RSS}{\operatorname{RSS}}
\newcommand{\SSE}{\operatorname{SSE}}
\newcommand{\SSR}{\operatorname{SSR}}
\newcommand{\logit}{\operatorname{logit}}
\newcommand{\se}{\operatorname{se}}
\newcommand{\SE}{\operatorname{SE}}
\newcommand{\diag}{\operatorname{diag}}
\newcommand{\R}{\mathbb{R}}
\newcommand{\answer}[1]{\boxed{\displaystyle #1}}

\newenvironment{resultbox}
  {\par\medskip\noindent\begin{minipage}{\linewidth}
   \color{navy!90!black}\hrule\medskip\textbf{Result.}\ }
  {\medskip\hrule\end{minipage}\par\medskip}

\newenvironment{checkbox}
  {\par\medskip\noindent\begin{minipage}{\linewidth}
   \color{navy!90!black}\hrule\medskip
   \textbf{\sffamily Calculator verification (Casio fx-580VN X).}\ }
  {\medskip\hrule\end{minipage}\par\medskip}

\begin{document}
\hypersetup{pageanchor=false}

\begin{titlepage}
  \centering
  \vspace*{0.85in}
  {\sffamily\color{navy}\Huge\bfseries STAT452\par}
  \vspace{0.20in}
  {\sffamily\color{teal}\LARGE Take-home Mid-term Retake\par}
  \vspace{0.15in}
  {\sffamily\color{navy!85}\LARGE Comprehensive Solutions\par}
  \vspace{0.48in}
  {\color{navy!40}\rule{0.78\textwidth}{1.1pt}\par}
  \vspace{0.48in}
  {\Large Applied Statistics for Engineers and Scientists II\par}
  \vspace{0.16in}
  {\large Term 3, Academic year 2025--2026 \(\cdot\) Class 24A02\par}
  \vfill
  \begin{minipage}{0.78\textwidth}
    \colorbox{softblue}{\parbox{0.94\linewidth}{
      \textbf{Conventions.} Natural logarithms are used throughout.
      All tests and intervals use \(\alpha_{\mathrm{sig}}=0.05\). Notation follows the
      Chapter 1--4 course notes, with exam-specific symbols identified
      explicitly. Intermediate calculations retain full precision; displayed
      decimal answers are rounded.
    }}
  \end{minipage}
  \vfill
  \begin{tabular}{@{}p{0.47\textwidth}p{0.35\textwidth}@{}}
    \textbf{Student name:} \hrulefill &
    \textbf{Student ID:} \hrulefill
  \end{tabular}
  \vspace{0.55in}
\end{titlepage}

\hypersetup{pageanchor=true}
\pagenumbering{roman}
\setcounter{tocdepth}{1}
\tableofcontents
\clearpage
\pagenumbering{arabic}

\section*{Notation crosswalk}
\addcontentsline{toc}{section}{Notation crosswalk}

The following conventions match the Chapter 1--4 notes and the exam.
\begin{itemize}
\item In Problems 1--2, \(p\) is the number of predictors \emph{excluding}
  the intercept. Thus \(X\) has \(p+1\) columns and
  \(df_E=n-p-1\). In Problem 1, \(p=2\) and \(df_E=7\).
  Problem 4 follows its own statement: there, \(X\in\mathbb R^{n\times p}\)
  and \(p\) is the number of columns of \(X\), equivalently the dimension of
  \(\bm\beta\).
\item Uppercase symbols such as \(Y_i\), \(Y\), and \(X\) denote random
  variables or model quantities; lowercase \(y_i\), \(x_i\), and \(v_i\)
  denote observed data. Bold symbols denote vectors.
\item The notes use
  \(\mathrm{SST}=\sum_i(y_i-\bar y)^2\),
  \(\mathrm{SSE}=\sum_i(y_i-\hat y_i)^2\), and
  \(\mathrm{SSR}=\sum_i(\hat y_i-\bar y)^2\).
  The exam's \(\mathrm{RSS}\) is the same quantity as the notes'
  \(\mathrm{SSE}\).
\item The significance level is written \(\alpha_{\mathrm{sig}}=0.05\) in
  this solution. The exam separately uses \(\alpha\) for the centered
  intercept in Problem 1 and for the positive scale parameter in Problem 2.
\item In Problem 2, \(a=\log\alpha\) is the log-scale intercept and
  corresponds to \(\beta_0\) in the Chapter 2 semi-log model.
\item Chapter 4 writes \(\pi(\bm X)=P(Y=1\mid\bm X)\) and uses threshold
  \(\tau\). Problem 3 uses \(\hat p(\bm x)\) and \(c\), so
  \(\hat p(\bm x)=\hat\pi(\bm x)\) and \(c=\tau\).
\item \(L(\alpha,\beta_1,\beta_2)\) is the least-squares criterion in Problem 1,
  matching Chapter 1. The exam independently defines \(L(c)\) as the
  threshold cost in Problem 3; the arguments distinguish the two functions.
\end{itemize}

\clearpage
\section*{Quick answer sheet}
\addcontentsline{toc}{section}{Quick answer sheet}

\begin{longtable}{>{\raggedright\arraybackslash}p{0.15\textwidth}
                  >{\raggedright\arraybackslash}p{0.77\textwidth}}
\toprule
\rowcolor{navy!10}\textbf{Item} & \textbf{Main result}\\
\midrule
\endfirsthead
\toprule
\rowcolor{navy!10}\textbf{Item} & \textbf{Main result}\\
\midrule
\endhead
1(a) & \(\hat\beta_0=7,\ \hat\beta_1=1.5,\ \hat\beta_2=-1\).\\
1(b) & \(\SSR=35,\ \SSE=5,\ R^2=0.875,\ R_{\rm adj}^2=47/56=0.8393,\ 
          \hat\sigma^2=5/7\).\\
1(c) & \(\widehat{\Var}(\hat{\bm\beta})
 =\begin{psmallmatrix}
 29/42&-5/42&-1/21\\[-1mm]
 -5/42&1/21&-1/42\\[-1mm]
 -1/21&-1/42&1/21
 \end{psmallmatrix}\).\\
1(d) & \(t=-\sqrt{21}=-4.5826,\ F=21=t^2,\ p=0.002536\); reject \(H_0\).\\
1(e) & \(\hat y_0=12.5\); mean CI \((11.405,13.595)\);
        prediction interval \((10.221,14.779)\).\\
2(a) & \(\hat a=1,\ \hat\beta=0.5,\ \hat\alpha=e=2.7183\);
        \(e^{\hat\beta}=1.6487\).\\
2(b) & \(\hat\sigma^2=0.025,\ \se(\hat\beta)=0.03780\);
        CI for \(\beta\): \((0.3951,0.6049)\);
        CI for \(e^\beta\): \((1.4845,1.8311)\).\\
2(c) & Estimated median \(=e^4=54.5982\);
        mean-corrected estimate \(=e^{4.0125}=55.2849\).\\
2(d) & Log-scale PI \((3.4003,4.5997)\);
        response-scale PI \((29.9736,99.4529)\).\\
3(a)--(b) & Conditional odds ratios:
        \(1.0725,\ 3.0042,\ 0.2725\).
        \(A:\hat p=0.1545\Rightarrow 0\);
        \(B:\hat p=0.8520\Rightarrow 1\).\\
3(c)--(d) & At \(c=0.50\): all five metrics equal \(0.50\) and
        \(L(0.50)=10\). At \(c=0.30\): accuracy \(0.50\), precision \(0.50\),
        recall \(0.75\), specificity \(0.25\), \(F_1=0.60\), and
        \(L(0.30)=7\). Select \(c=0.30\).\\
4 & For \(D=C-(X^\trans X)^{-1}X^\trans\),
    \(\Var(\tilde{\bm\beta})-\Var(\hat{\bm\beta})
      =\sigma^2DD^\trans\succeq0\).\\
\bottomrule
\end{longtable}

\clearpage
\section{Problem 1: Multiple linear regression}

The model is
\[
  Y_i=\beta_0+\beta_1x_{i1}+\beta_2x_{i2}+\varepsilon_i,
  \qquad \varepsilon_i\stackrel{\mathrm{iid}}{\sim}N(0,\sigma^2),
\]
with \(n=10\). Define
\[
  \bm z_i=
  \begin{pmatrix}x_{i1}-4\\x_{i2}-3\end{pmatrix},
  \qquad w_i=y_i-10,
\]
so that \(\sum_i\bm z_i=\bm0\) and \(\sum_iw_i=0\). The given cross-products are
\[
  G=\sum_{i=1}^{10}\bm z_i\bm z_i^\trans
   =\begin{pmatrix}20&10\\10&20\end{pmatrix},
  \qquad
  \bm g=\sum_{i=1}^{10}\bm z_iw_i
   =\begin{pmatrix}20\\-5\end{pmatrix}.
\]

\subsection{Part (a): Centered normal equations and estimates}

Write \(z_{i1}=x_{i1}-4\) and \(z_{i2}=x_{i2}-3\), so that
\(\bm z_i=(z_{i1},z_{i2})^\trans\). The centered intercept specified by the
exam is
\[
  \alpha=\beta_0+4\beta_1+3\beta_2.
\]
Hence
\[
\beta_0+\beta_1x_{i1}+\beta_2x_{i2}
=\alpha+\beta_1z_{i1}+\beta_2z_{i2}.
\]
Because \(y_i=\bar y+w_i\), the least-squares criterion becomes
\begin{align*}
L(\alpha,\beta_1,\beta_2)
 &=\sum_{i=1}^{10}
 \left[y_i-\{\alpha+\beta_1z_{i1}+\beta_2z_{i2}\}\right]^2\\
 &=\sum_{i=1}^{10}
 \left[w_i-\beta_1z_{i1}-\beta_2z_{i2}-(\alpha-\bar y)\right]^2.
\end{align*}
Differentiating with respect to \(\alpha\), and using
\(\sum_iw_i=\sum_i z_{i1}=\sum_i z_{i2}=0\),
\[
 \frac{\partial L}{\partial\alpha}
 =-2\sum_{i=1}^{10}
 \left[w_i-\beta_1z_{i1}-\beta_2z_{i2}-(\alpha-\bar y)\right]
 =20(\alpha-\bar y).
\]
Thus \(\hat\alpha=\bar y=10\). At this value,
\[
 L(10,\beta_1,\beta_2)
 =\sum_{i=1}^{10}(w_i-\beta_1z_{i1}-\beta_2z_{i2})^2.
\]
For \(j=1,2\), the slope normal equations are
\[
\frac{\partial L}{\partial\beta_j}
=-2\sum_{i=1}^{10}z_{ij}
 (w_i-\beta_1z_{i1}-\beta_2z_{i2})=0.
\]
Stacking these two equations gives
\[
\begin{pmatrix}
\sum_i z_{i1}^2&\sum_i z_{i1}z_{i2}\\
\sum_i z_{i1}z_{i2}&\sum_i z_{i2}^2
\end{pmatrix}
\begin{pmatrix}\hat\beta_1\\\hat\beta_2\end{pmatrix}
=
\begin{pmatrix}\sum_i z_{i1}w_i\\\sum_i z_{i2}w_i\end{pmatrix},
\]
which, using the notation supplied by the exam, is
\[
  \answer{G
  \begin{pmatrix}\hat\beta_1\\\hat\beta_2\end{pmatrix}=\bm g}.
\]
Now
\[
G^{-1}
=\frac{1}{20(20)-10(10)}
  \begin{pmatrix}20&-10\\-10&20\end{pmatrix}
=\begin{pmatrix}
1/15&-1/30\\
-1/30&1/15
\end{pmatrix}.
\]
Therefore,
\[
\begin{pmatrix}\hat\beta_1\\\hat\beta_2\end{pmatrix}
=G^{-1}\bm g
=\begin{pmatrix}
1/15&-1/30\\
-1/30&1/15
\end{pmatrix}
\begin{pmatrix}20\\-5\end{pmatrix}
=\begin{pmatrix}3/2\\-1\end{pmatrix}.
\]
Finally,
\[
\hat\beta_0
=\hat\alpha-4\hat\beta_1-3\hat\beta_2
=10-4(3/2)-3(-1)=7.
\]

\begin{resultbox}
\[
  \answer{\hat y=7+1.5x_1-x_2}.
\]
\end{resultbox}

\subsection{Part (b): Regression sums of squares and variance estimate}

In centered regression, the explained sum of squares is
\[
 \SSR=\begin{pmatrix}\hat\beta_1&\hat\beta_2\end{pmatrix}\bm g
 =\begin{pmatrix}3/2&-1\end{pmatrix}
  \begin{pmatrix}20\\-5\end{pmatrix}
 =30+5=35.
\]
Since \(\mathrm{SST}=40\),
\[
 \SSE=\mathrm{SST}-\SSR=40-35=5.
\]
Thus
\[
 R^2=\frac{\SSR}{\mathrm{SST}}=\frac{35}{40}
 =\answer{0.875}.
\]
There are \(p=2\) predictors and \(p+1=3\) fitted coefficients, so the residual
degrees of freedom are \(df_E=n-p-1=10-2-1=7\). Hence
\begin{align*}
R_{\mathrm{adj}}^2
 &=1-\frac{\SSE/(n-p-1)}{\mathrm{SST}/(n-1)}
 =1-\frac{5/7}{40/9}
 =\frac{47}{56}
 =\answer{0.839286},\\
\hat\sigma^2=\mathrm{MSE}
&=\frac{\SSE}{n-p-1}=\answer{\frac57=0.714286},
\qquad
\hat\sigma=\sqrt{\frac57}=\answer{0.845154}.
\end{align*}

\subsection{Part (c): Covariance under centering and transformation}

Let \(Z\) be the \(10\times2\) matrix whose \(i\)-th row is
\(\bm z_i^\trans\), and let the centered design matrix be
\[
 W=\begin{pmatrix}\bm1&Z\end{pmatrix}.
\]
Because every column of \(Z\) is centered,
\(\bm1^\trans Z=\bm0^\trans\), while \(Z^\trans Z=G\). Therefore
\[
 W^\trans W
 =\begin{pmatrix}
 \bm1^\trans\bm1&\bm1^\trans Z\\
 Z^\trans\bm1&Z^\trans Z
 \end{pmatrix}
 =\begin{pmatrix}10&\bm0^\trans\\\bm0&G\end{pmatrix}.
\]
For \(\bm\theta=(\alpha,\beta_1,\beta_2)^\trans\), the OLS covariance is
\[
\Var(\hat{\bm\theta})
=\sigma^2(W^\trans W)^{-1}
=\sigma^2
\begin{pmatrix}
1/10&\bm0^\trans\\
\bm0&G^{-1}
\end{pmatrix}.
\]
Replacing \(\sigma^2\) by \(5/7\) gives the estimated version requested in
the question.

To return to the original parameterization, use
\[
\begin{pmatrix}\beta_0\\\beta_1\\\beta_2\end{pmatrix}
=A\begin{pmatrix}\alpha\\\beta_1\\\beta_2\end{pmatrix},
\qquad
A=
\begin{pmatrix}
1&-4&-3\\
0&1&0\\
0&0&1
\end{pmatrix}.
\]
Thus
\begin{align*}
\widehat{\Var}(\hat{\bm\beta})
&=\hat\sigma^2 A
\begin{pmatrix}
1/10&0&0\\
0&1/15&-1/30\\
0&-1/30&1/15
\end{pmatrix}A^\trans\\
&=\frac57
\begin{pmatrix}
29/30&-1/6&-1/15\\
-1/6&1/15&-1/30\\
-1/15&-1/30&1/15
\end{pmatrix}\\
&=\answer{
\begin{pmatrix}
29/42&-5/42&-1/21\\
-5/42&1/21&-1/42\\
-1/21&-1/42&1/21
\end{pmatrix}}.
\end{align*}
In decimal form, this is
\[
\begin{pmatrix}
0.690476&-0.119048&-0.047619\\
-0.119048&0.047619&-0.023810\\
-0.047619&-0.023810&0.047619
\end{pmatrix}.
\]
Notice why centering is useful: \(\hat\alpha\) is uncorrelated with both
slope estimators, even though \(\hat\beta_0\) generally is not.

\subsection{Part (d): Testing \(H_0:\beta_2=0\)}

\subsubsection{Coefficient \(t\) test}

From the \((3,3)\) entry of the covariance matrix,
\[
\SE(\hat\beta_2)
=\sqrt{\hat\sigma^2(G^{-1})_{22}}
=\sqrt{\frac57\cdot\frac1{15}}
=\sqrt{\frac1{21}}.
\]
Therefore,
\[
t=\frac{\hat\beta_2-0}{\SE(\hat\beta_2)}
=\frac{-1}{\sqrt{1/21}}
=-\sqrt{21}
=\answer{-4.582576},
\qquad df=7.
\]
The two-sided \(p\)-value is
\[
 2P(T_7\geq |t|)=\answer{0.002536}.
\]

\subsubsection{Nested-model \(F\) test}

The reduced model has \(\RSS_{\mathrm{red}}=20\), while the full model has
\(\RSS_{\mathrm{full}}=\SSE=5\). There is \(q=1\) restriction, so
\[
F
=\frac{(\RSS_{\mathrm{red}}-\RSS_{\mathrm{full}})/q}
       {\RSS_{\mathrm{full}}/(n-p-1)}
=\frac{(20-5)/1}{5/7}
=\answer{21},
\qquad (df_1,df_2)=(1,7).
\]
Also,
\[
 t^2=(-\sqrt{21})^2=21=F.
\]
Both methods therefore give the same \(p\)-value and the same decision:
\[
\answer{\text{Reject \(H_0\) at the 5\% level.}}
\]
After controlling for \(x_1\), the data provide strong evidence that the
coefficient of \(x_2\) is nonzero (and its fitted sign is negative).

\subsection{Part (e): Mean-response and prediction intervals at \(x_0=(1,5,2)^\trans\)}

The centered predictor vector is
\[
\bm z_0=
\begin{pmatrix}5-4\\2-3\end{pmatrix}
=\begin{pmatrix}1\\-1\end{pmatrix}.
\]
The fitted response is
\[
\hat y_0=7+1.5(5)-1(2)=\answer{12.5}.
\]
The leverage is
\begin{align*}
h_0
&=\frac1{10}+\bm z_0^\trans G^{-1}\bm z_0\\
&=\frac1{10}
+\begin{pmatrix}1&-1\end{pmatrix}
\begin{pmatrix}1/15&-1/30\\-1/30&1/15\end{pmatrix}
\begin{pmatrix}1\\-1\end{pmatrix}\\
&=\frac1{10}+\frac15=\answer{\frac3{10}}.
\end{align*}

For the conditional mean, the standard error is
\[
\hat\sigma\sqrt{h_0}
=\sqrt{\frac57\cdot\frac3{10}}
=\sqrt{\frac3{14}}
=0.462910.
\]
Using \(t_{0.975,7}=2.365\),
\[
12.5\pm2.365(0.462910)
=12.5\pm1.094608,
\]
so the 95\% confidence interval for \(\E(Y\mid x_0)\) is
\[
\answer{(11.4054,\ 13.5946)}.
\]

For a new response, the standard error includes the new error term:
\[
\hat\sigma\sqrt{1+h_0}
=\sqrt{\frac57\left(1+\frac3{10}\right)}
=\sqrt{\frac{13}{14}}
=0.963624.
\]
Thus the 95\% prediction interval is
\[
12.5\pm2.365(0.963624)
=12.5\pm2.278609
=\answer{(10.2214,\ 14.7786)}.
\]
The prediction interval is wider because it accounts for both uncertainty
in the estimated mean and irreducible variation in a future response.

\begin{checkbox}
\begin{enumerate}
\item In \texttt{MENU \(\to\) Matrix}, define
\[
\mathrm{MatA}=\begin{pmatrix}20&10\\10&20\end{pmatrix},
\qquad
\mathrm{MatB}=\begin{pmatrix}20\\-5\end{pmatrix}.
\]
Evaluate \(\mathrm{MatA}^{-1}\mathrm{MatB}\). The display should be
\((1.5,-1)^\trans\).
\item In ordinary Calculate mode, enter
\(\sqrt{(5\div7)(1\div15)}\). The display should be \(0.2182179\);
then \(-1\div\text{Ans}\) gives \(-4.582576\).
\item Enter \((20-5)\div(5\div7)\). The display should be \(21\), confirming
\(F=t^2\).
\item For part (e), enter
\(12.5\pm2.365\sqrt{(5\div7)(3\div10)}\), then replace \(3\div10\)
inside the square root by \(1+3\div10\). This reproduces the CI and PI
endpoints above.
\end{enumerate}
\end{checkbox}

\clearpage
\section{Problem 2: Transformation and back-transformation}

The response model is
\[
Y=\alpha\exp(\beta X)\exp(\varepsilon),
\qquad \varepsilon\sim N(0,\sigma^2),\quad \alpha>0.
\]

\subsection{Part (a): Linearization, estimation, and interpretation}

Taking natural logarithms and writing \(V=\log Y\) gives
\[
V=\log\alpha+\beta X+\varepsilon=a+\beta X+\varepsilon,
\qquad a=\log\alpha.
\]
Thus \(a\) is exactly the log-scale intercept denoted by \(\beta_0\) in the
Chapter 2 semi-log model; the symbol \(a\) is retained here because it is the
notation specified by the exam. This is an ordinary simple linear regression
on the log-response scale.
The least-squares slope is
\[
\hat\beta=\frac{S_{xv}}{S_{xx}}
=\frac{35/4}{35/2}=\answer{\frac12}.
\]
The intercept is
\[
\hat a=\bar v-\hat\beta\bar x
=\frac94-\frac12\left(\frac52\right)
=\frac94-\frac54=\answer{1}.
\]
Therefore,
\[
\hat\alpha=\exp(\hat a)=\answer{e=2.718282}.
\]

Holding all else fixed, increasing \(X\) by one changes the fitted median from
\(\exp(\hat a+\hat\beta x)\) to
\(\exp(\hat a+\hat\beta(x+1))\). Their ratio is
\[
\exp(\hat\beta)=e^{1/2}=\answer{1.648721}.
\]
Thus each one-unit increase in \(X\) multiplies the conditional median of
\(Y\) by about \(1.6487\), corresponding to an estimated \(64.87\%\) increase.

\subsection{Part (b): Error variance and confidence interval for the effect}

There are \(n-2=4\) residual degrees of freedom. Hence
\[
\hat\sigma^2=\mathrm{MSE}_{\log}=\frac{\SSE_{\log}}{n-2}
=\frac{0.10}{4}=\answer{0.025},
\qquad
\hat\sigma=\sqrt{0.025}=0.158114.
\]
The slope standard error is
\[
\se(\hat\beta)
=\sqrt{\frac{\hat\sigma^2}{S_{xx}}}
=\sqrt{\frac{0.025}{35/2}}
=\sqrt{\frac1{700}}
=\answer{0.037796}.
\]
Using \(t_{0.975,4}=2.776\),
\begin{align*}
\hat\beta\pm t_{0.975,4}\se(\hat\beta)
&=0.5\pm2.776(0.037796)\\
&=0.5\pm0.104923.
\end{align*}
Thus the 95\% confidence interval for \(\beta\) is
\[
\answer{(0.395077,\ 0.604923)}.
\]
Since the exponential function is strictly increasing, exponentiating the
two endpoints gives the 95\% interval for the one-unit multiplicative effect:
\[
\answer{(e^{0.395077},e^{0.604923})
=(1.484499,\ 1.831111)}.
\]
In words, a one-unit increase in \(X\) is associated with a multiplication
of the conditional median by between about \(1.48\) and \(1.83\).

\subsection{Part (c): Conditional median and lognormal mean correction}

At \(x_0=6\),
\[
\hat v_0=\hat a+\hat\beta x_0=1+\frac12(6)=4.
\]
Because a normal error has median zero, the fitted conditional median is
\[
\widehat{\operatorname{Med}}(Y\mid X=6)
=\exp(\hat v_0)=\answer{e^4=54.598150}.
\]
The stated lognormal mean correction gives
\[
\widehat{\E}(Y\mid X=6)
=\exp\left(\hat v_0+\frac{\hat\sigma^2}{2}\right)
=\exp\left(4+\frac{0.025}{2}\right)
=\answer{e^{4.0125}=55.284913}.
\]
They differ because, although \(\operatorname{Med}(e^\varepsilon)=1\),
\[
\E(e^\varepsilon)=e^{\sigma^2/2}>1.
\]
Exponentiation makes the response distribution right-skewed; large positive
errors increase the mean more than equally sized negative errors decrease it.
The mean is therefore larger than the median.

\subsection{Part (d): Prediction interval at \(x_0=6\)}

The leverage for predicting on the log scale is
\begin{align*}
h_0
&=\frac16+\frac{(x_0-\bar x)^2}{S_{xx}}\\
&=\frac16+\frac{(6-5/2)^2}{35/2}
=\frac16+\frac{(7/2)^2}{35/2}\\
&=\frac16+\frac7{10}
=\answer{\frac{13}{15}=0.866667}.
\end{align*}
The prediction standard error on the log scale is
\[
\hat\sigma\sqrt{1+h_0}
=\sqrt{0.025\left(1+\frac{13}{15}\right)}
=\sqrt{\frac7{150}}
=0.216025.
\]
Thus
\[
\hat v_0\pm t_{0.975,4}\hat\sigma\sqrt{1+h_0}
=4\pm2.776(0.216025)
=4\pm0.599685.
\]
The 95\% prediction interval for \(V_0=\log Y_0\) is
\[
\answer{(3.400315,\ 4.599685)}.
\]
Exponentiating both endpoints,
\[
\answer{(e^{3.400315},e^{4.599685})
=(29.973550,\ 99.452943)}
\]
is the 95\% prediction interval for a future response \(Y_0\). It is
asymmetric about \(54.598\) on the original scale, as expected after
exponentiation.

\begin{checkbox}
\begin{enumerate}
\item In Calculate mode, enter
\((35\div4)\div(35\div2)\); the display should be \(0.5\).
Then enter \(9\div4-0.5(5\div2)\); the display should be \(1\).
\item Enter
\(\sqrt{0.025\div(35\div2)}\); the display should be \(0.03779645\).
The two expressions
\(0.5-2.776\,\text{Ans}\) and \(0.5+2.776\,\text{Ans}\)
give the slope interval endpoints.
\item Use the \(e^x\) function on those endpoints to obtain
\(1.484499\) and \(1.831111\).
\item Enter \(e^4\) and \(e^{4+0.025/2}\) to obtain \(54.59815\) and
\(55.28491\). For the prediction interval, evaluate
\(e^{\,4\pm2.776\sqrt{0.025(1+13/15)}}\).
\end{enumerate}
\end{checkbox}

\clearpage
\section{Problem 3: Logistic regression evaluation}

The fitted logit is
\[
\log\left(\frac{\hat p(\bm x)}{1-\hat p(\bm x)}\right)
=-5+0.07x_1+1.10x_2-1.30x_3.
\]
In the notation of Chapter 4,
\(\pi(\bm X)=P(Y=1\mid\bm X)\), so
\(\hat p(\bm x)=\hat\pi(\bm x)\). The exam's cutoff \(c\) is the same
classification threshold denoted by \(\tau\) in the notes.

Following Chapter 4 directly, the fitted probability is obtained by applying
its logistic function to the fitted log-odds:
\[
\hat p(\bm x)
=\frac{1}{1+\exp\!\left[-\{-5+0.07x_1+1.10x_2-1.30x_3\}\right]}.
\]

\subsection{Part (a): Conditional odds-ratio interpretations}

Each interpretation holds with the other predictors fixed.
\begin{itemize}
\item A one-unit increase in temperature \(x_1\) multiplies the alarm odds by
\[
e^{0.07}=\answer{1.072508},
\]
an increase of approximately \(7.25\%\).
\item A one-unit increase in vibration \(x_2\) multiplies the alarm odds by
\[
e^{1.10}=\answer{3.004166},
\]
so the odds are approximately tripled.
\item Changing from no preventive maintenance (\(x_3=0\)) to preventive
maintenance (\(x_3=1\)) multiplies the alarm odds by
\[
e^{-1.30}=\answer{0.272532}.
\]
Thus preventive maintenance is associated with a \(1-0.272532=72.75\%\)
reduction in the conditional odds of an alarm. Equivalently, the odds without
maintenance are \(e^{1.30}=3.6693\) times those with maintenance.
\end{itemize}

\subsection{Part (b): Probabilities and classification at threshold \(0.5\)}

For \(A=(50,1,1)\),
\[
\log\left(\frac{\hat p_A}{1-\hat p_A}\right)
=-5+0.07(50)+1.10(1)-1.30(1)=-1.70,
\]
and therefore
\[
\hat p_A=\frac1{1+e^{1.70}}=\answer{0.154465}.
\]
Since \(\hat p_A<0.5\), the fitted class is \(\answer{0}\).

For \(B=(65,2,0)\),
\[
\log\left(\frac{\hat p_B}{1-\hat p_B}\right)
=-5+0.07(65)+1.10(2)-1.30(0)=1.75,
\]
so
\[
\hat p_B=\frac1{1+e^{-1.75}}=\answer{0.851953}.
\]
Since \(\hat p_B\geq0.5\), the fitted class is \(\answer{1}\).

\subsection{Part (c): Confusion matrices and metrics}

The metric definitions are
\[
\begin{aligned}
\text{accuracy}&=\frac{TP+TN}{TP+FP+TN+FN},&
\text{precision}&=\frac{TP}{TP+FP},\\
\text{recall}&=\frac{TP}{TP+FN},&
\text{specificity}&=\frac{TN}{TN+FP},\\
F_1&=\frac{2(\text{precision})(\text{recall})}
{\text{precision}+\text{recall}}
=\frac{2TP}{2TP+FP+FN}.
\end{aligned}
\]

\subsubsection{Threshold \(c=0.50\)}

The predicted labels are
\[
(\hat Y_1,\ldots,\hat Y_8)=(1,1,1,1,0,0,0,0).
\]
Comparing these to \((1,0,1,0,1,0,1,0)\) gives
\[
TP=2,\qquad FP=2,\qquad TN=2,\qquad FN=2.
\]
Hence
\[
\begin{aligned}
\text{accuracy}&=\frac{2+2}{8}=0.50,&
\text{precision}&=\frac{2}{2+2}=0.50,\\
\text{recall}&=\frac{2}{2+2}=0.50,&
\text{specificity}&=\frac{2}{2+2}=0.50,\\
F_1&=\frac{2(2)}{2(2)+2+2}=\answer{0.50}.
\end{aligned}
\]

\subsubsection{Threshold \(c=0.30\)}

The predicted labels are
\[
(\hat Y_1,\ldots,\hat Y_8)=(1,1,1,1,1,1,0,0).
\]
Therefore,
\[
TP=3,\qquad FP=3,\qquad TN=1,\qquad FN=1.
\]
The metrics are
\[
\begin{aligned}
\text{accuracy}&=\frac{3+1}{8}=\answer{0.50},&
\text{precision}&=\frac{3}{3+3}=\answer{0.50},\\
\text{recall}&=\frac{3}{3+1}=\answer{0.75},&
\text{specificity}&=\frac{1}{1+3}=\answer{0.25},\\
F_1&=\frac{2(3)}{2(3)+3+1}=\answer{0.60}.
\end{aligned}
\]

\begin{center}
\begin{tabular}{@{}lrrrrrrr@{}}
\toprule
\rowcolor{navy!10}
\(c\)&\(TP\)&\(FP\)&\(TN\)&\(FN\)&Accuracy&Precision&Recall\\
\midrule
0.50&2&2&2&2&0.50&0.50&0.50\\
0.30&3&3&1&1&0.50&0.50&0.75\\
\bottomrule
\end{tabular}

\vspace{0.5em}
\begin{tabular}{@{}lrr@{}}
\toprule
\rowcolor{navy!10}
\(c\)&Specificity&\(F_1\)\\
\midrule
0.50&0.50&0.50\\
0.30&0.25&0.60\\
\bottomrule
\end{tabular}
\end{center}

Lowering the threshold finds one additional true alarm (recall rises), but it
also creates one additional false alarm (specificity falls). Accuracy and
precision happen to remain unchanged in this test set.

\subsection{Part (d): Cost-sensitive threshold choice}

Using \(L(c)=4FN(c)+FP(c)\),
\[
L(0.50)=4(2)+2=\answer{10},
\qquad
L(0.30)=4(1)+3=\answer{7}.
\]
Therefore,
\[
\answer{\text{Select \(c=0.30\).}}
\]
Although both thresholds have the same accuracy, the lower threshold avoids
one costly false negative. Its one extra false positive costs only \(1\), so
the total empirical cost decreases by \(3\).

\begin{checkbox}
In Calculate mode:
\begin{enumerate}
\item For case \(A\), enter
\(-5+0.07(50)+1.10(1)-1.30(1)\); the display is \(-1.7\).
Then enter \(1\div(1+e^{-\text{Ans}})\); the display is \(0.1544653\).
\item Repeat with \(B\); the linear predictor is \(1.75\) and the probability
is \(0.8519528\).
\item For each threshold, write \(1\) beside every probability at least as
large as \(c\), then tally the four cells. The two tallies should be
\((TP,FP,TN,FN)=(2,2,2,2)\) and \((3,3,1,1)\).
\item Enter \(4(2)+2\) and \(4(1)+3\); the displays \(10\) and \(7\) verify
the cost comparison.
\end{enumerate}
\end{checkbox}

\clearpage
\section{Problem 4: Gauss--Markov theorem}

Assume
\[
\bm Y=X\bm\beta+\bm\varepsilon,
\qquad
\E(\bm\varepsilon)=\bm0,
\qquad
\Var(\bm\varepsilon)=\sigma^2I_n,
\]
where \(X\in\R^{n\times p}\) is deterministic and
\(\operatorname{rank}(X)=p\). Here \(p\) follows the notation of Problem 4:
it is the number of columns of \(X\) and the dimension of \(\bm\beta\), not
the predictor count convention used in Problem 1. No distributional assumption
beyond these first two moments is needed.

\subsection{Step 1: Error and covariance of OLS}

Let
\[
A=(X^\trans X)^{-1}X^\trans,
\qquad
\hat{\bm\beta}=A\bm Y.
\]
Full column rank guarantees that \(X^\trans X\) is invertible, and
\[
AX=(X^\trans X)^{-1}X^\trans X=I_p.
\]
Substituting \(\bm Y=X\bm\beta+\bm\varepsilon\),
\begin{align*}
\hat{\bm\beta}
&=(X^\trans X)^{-1}X^\trans(X\bm\beta+\bm\varepsilon)\\
&=\bm\beta+(X^\trans X)^{-1}X^\trans\bm\varepsilon.
\end{align*}
Therefore,
\[
\answer{\hat{\bm\beta}-\bm\beta
=(X^\trans X)^{-1}X^\trans\bm\varepsilon}.
\]
Its covariance is
\begin{align*}
\Var(\hat{\bm\beta})
&=A\Var(\bm\varepsilon)A^\trans\\
&=\sigma^2
(X^\trans X)^{-1}X^\trans X(X^\trans X)^{-1}\\
&=\answer{\sigma^2(X^\trans X)^{-1}}.
\end{align*}

\subsection{Step 2: What unbiasedness of a general linear estimator implies}

Let \(\tilde{\bm\beta}=C\bm Y\), where \(C\in\R^{p\times n}\) is deterministic.
Then
\[
\E(\tilde{\bm\beta})
=C\E(\bm Y)
=CX\bm\beta.
\]
Unbiasedness for every \(\bm\beta\in\R^p\) requires
\[
CX\bm\beta=\bm\beta
\quad\text{for every }\bm\beta.
\]
Equivalently, \((CX-I_p)\bm\beta=\bm0\) for every \(\bm\beta\), which can
hold only when
\[
\answer{CX=I_p}.
\]

\subsection{Step 3: Decomposition and vanishing cross terms}

Define
\[
D=C-A=C-(X^\trans X)^{-1}X^\trans.
\]
Using \(CX=I_p\) and \(AX=I_p\),
\[
DX=CX-AX=I_p-I_p=\answer{0}.
\]
Since \(C=A+D\),
\begin{align*}
\Var(\tilde{\bm\beta})
&=C\Var(\bm Y)C^\trans\\
&=\sigma^2(A+D)(A+D)^\trans\\
&=\sigma^2\left(AA^\trans+AD^\trans+DA^\trans+DD^\trans\right).
\end{align*}
The cross terms are zero. Indeed,
\[
AD^\trans=(X^\trans X)^{-1}X^\trans D^\trans
=(X^\trans X)^{-1}(DX)^\trans=0,
\]
and, equivalently by transposition,
\[
DA^\trans=DX(X^\trans X)^{-1}=0.
\]
Moreover, \(\Var(\hat{\bm\beta})=\sigma^2AA^\trans\). Consequently,
\[
\answer{
\Var(\tilde{\bm\beta})-\Var(\hat{\bm\beta})
=\sigma^2DD^\trans}.
\]
For every \(\bm u\in\R^p\),
\[
\bm u^\trans(\sigma^2DD^\trans)\bm u
=\sigma^2(D^\trans\bm u)^\trans(D^\trans\bm u)
=\sigma^2\lVert D^\trans\bm u\rVert^2\geq0.
\]
Hence \(\sigma^2DD^\trans\) is positive semidefinite, proving
\[
\boxed{
\Var(\tilde{\bm\beta})-\Var(\hat{\bm\beta})\succeq0}.
\]

\begin{resultbox}
Among all linear estimators that are unbiased for every \(\bm\beta\), OLS has
the smallest covariance matrix in the positive-semidefinite (Loewner) order.
The proof used only linearity, unbiasedness, and
\(\Var(\bm\varepsilon)=\sigma^2I_n\); it did not use Gaussian errors.
\end{resultbox}

\begin{checkbox}
No numerical calculator step is needed for this proof. A useful dimension
check is
\[
C,D,A\in\R^{p\times n},\quad X\in\R^{n\times p},\quad
DX\in\R^{p\times p},\quad DD^\trans\in\R^{p\times p}.
\]
These dimensions confirm that every matrix product in the covariance
decomposition is conformable.
\end{checkbox}

\clearpage
\appendix
\section{Formula reference aligned with Chapters 1--4}

This appendix collects the formulas used in the solutions with the notation of
the course notes. In the Chapter 1--3 regression formulas, \(n\) is the sample
size and \(p\) is the number of predictors excluding the intercept. Problem 4
retains the exam's separate convention \(X\in\mathbb R^{n\times p}\), where
\(p=\dim(\bm\beta)\). Formulas explicitly supplied by the exam are labelled as
such rather than attributed to the notes. To avoid the exam's reuse of
\(\alpha\), the significance level is denoted by \(\alpha_{\mathrm{sig}}\).

\subsection{Chapter 1: Multiple linear regression}

\subsubsection*{Model, least squares, and estimation}
\[
\bm Y=X\bm\beta+\bm\varepsilon,
\qquad
\bm\varepsilon\sim N(\bm0,\sigma^2I_n),
\qquad
L(\bm\beta)=(\bm Y-X\bm\beta)^\trans(\bm Y-X\bm\beta).
\]
The normal equations and least-squares estimator are
\[
X^\trans X\hat{\bm\beta}=X^\trans\bm Y,
\qquad
\hat{\bm\beta}=(X^\trans X)^{-1}X^\trans\bm Y.
\]
With \(p\) predictors, the error degrees of freedom and variance estimate are
\[
df_E=n-p-1,
\qquad
\hat\sigma^2=\mathrm{MSE}
=\frac{\mathrm{SSE}}{n-p-1},
\qquad
\mathrm{SSE}=\sum_{i=1}^n(y_i-\hat y_i)^2.
\]

\subsubsection*{Covariance and coefficient inference}
\[
\Var(\hat{\bm\beta})=\sigma^2(X^\trans X)^{-1},
\qquad
\SE(\hat\beta_j)
=\sqrt{\mathrm{MSE}\,[(X^\trans X)^{-1}]_{jj}}.
\]
For \(H_0:\beta_j=\beta_{j,0}\),
\[
t=\frac{\hat\beta_j-\beta_{j,0}}{\SE(\hat\beta_j)},
\qquad t\sim t_{n-p-1}\quad\text{under }H_0,
\]
and the \(100(1-\alpha_{\mathrm{sig}})\%\) confidence interval is
\[
\hat\beta_j\pm t_{1-\alpha_{\mathrm{sig}}/2,n-p-1}\SE(\hat\beta_j).
\]

\subsubsection*{Sums of squares and goodness of fit}
\[
\mathrm{SST}=\sum_i(y_i-\bar y)^2,
\quad
\mathrm{SSR}=\sum_i(\hat y_i-\bar y)^2,
\quad
\mathrm{SSE}=\sum_i(y_i-\hat y_i)^2,
\quad
\mathrm{SST}=\mathrm{SSR}+\mathrm{SSE}.
\]
\[
R^2=\frac{\mathrm{SSR}}{\mathrm{SST}}
=1-\frac{\mathrm{SSE}}{\mathrm{SST}},
\qquad
R_{\mathrm{adj}}^2
=1-\frac{\mathrm{SSE}/(n-p-1)}{\mathrm{SST}/(n-1)}.
\]

\subsubsection*{Nested-model partial \(F\) test}
If \(q=df_{E,\mathrm{red}}-df_{E,\mathrm{full}}\) restrictions separate the
reduced and full models, then
\[
F=\frac{(\mathrm{SSE}_{\mathrm{red}}-
                 \mathrm{SSE}_{\mathrm{full}})/q}
        {\mathrm{MSE}_{\mathrm{full}}}
\sim F_{q,df_{E,\mathrm{full}}}\quad\text{under }H_0.
\]
The exam writes these residual sums of squares as \(\mathrm{RSS}\). For one
restriction, the coefficient test and nested-model test satisfy \(F=t^2\).

\subsubsection*{Prediction formulas supplied by the exam}
For a centered predictor vector \(\bm z_0\),
\[
h_0=\frac1n+\bm z_0^\trans G^{-1}\bm z_0.
\]
The mean-response confidence interval and new-response prediction interval are
\[
\hat y_0\pm t_{1-\alpha_{\mathrm{sig}}/2,n-p-1}\hat\sigma\sqrt{h_0},
\qquad
\hat y_0\pm t_{1-\alpha_{\mathrm{sig}}/2,n-p-1}\hat\sigma\sqrt{1+h_0}.
\]

\subsection{Chapter 2: Transformation of variables}

For the semi-log model in Problem 2,
\[
Y=\alpha e^{\beta X}e^\varepsilon
\quad\Longrightarrow\quad
V=\log Y=a+\beta X+\varepsilon,
\qquad a=\log\alpha.
\]
This is the Chapter 2 formula
\(\log Y=\beta_0+\beta_1X\), with \(a\equiv\beta_0\) and
\(\beta\equiv\beta_1\). The simple-regression estimators are
\[
\hat\beta=\frac{S_{xv}}{S_{xx}},
\qquad
\hat a=\bar v-\hat\beta\bar x,
\qquad
\hat\alpha=e^{\hat a}.
\]
A one-unit increase in \(X\) multiplies the conditional median of \(Y\) by
\[
e^{\hat\beta}.
\]
The slope standard error and interval are
\[
\se(\hat\beta)=\sqrt{\frac{\mathrm{MSE}_{\log}}{S_{xx}}},
\qquad
\hat\beta\pm t_{1-\alpha_{\mathrm{sig}}/2,n-2}\se(\hat\beta).
\]
Exponentiating the endpoints gives an interval for the multiplicative effect
\(e^\beta\).

For the lognormal model, the following back-transformations are supplied by
the exam:
\[
\widehat{\operatorname{Med}}(Y\mid X=x_0)
=e^{\hat a+\hat\beta x_0},
\qquad
\widehat{\E}(Y\mid X=x_0)
=e^{\hat a+\hat\beta x_0+\hat\sigma^2/2}.
\]
For a future log-response,
\[
h_0=\frac1n+\frac{(x_0-\bar x)^2}{S_{xx}},
\qquad
\hat v_0\pm t_{1-\alpha_{\mathrm{sig}}/2,n-2}\hat\sigma\sqrt{1+h_0}.
\]
Exponentiating both endpoints gives the prediction interval for \(Y_0\).

\subsection{Chapter 3: Relevance to this retake}

The retake does not ask for variable selection, cross-validation, ridge, or
LASSO calculations. The Chapter 3 Wald statistic
\[
t=\frac{\hat\beta_j}{\SE(\hat\beta_j)}
\]
is the same coefficient statistic used in Problem 1(d). Because Problem 1 is
an ordinary least-squares model with normally distributed errors, the exact
finite-sample reference distribution is \(t_{n-p-1}\), as stated in Chapter 1.

\subsection{Chapter 4: Logistic regression and classification}

The notes define
\[
\pi(\bm X)=P(Y=1\mid\bm X),
\qquad
\logit\{\pi(\bm X)\}
=\log\left(\frac{\pi(\bm X)}{1-\pi(\bm X)}\right)
=\bm X^\trans\bm\beta,
\]
and therefore
\[
\pi(\bm X)=\frac{e^{\bm X^\trans\bm\beta}}
{1+e^{\bm X^\trans\bm\beta}}
=\frac1{1+e^{-\bm X^\trans\bm\beta}}.
\]
The conditional odds and one-unit odds ratio are
\[
\operatorname{Odds}(Y=1\mid\bm X)=\frac{\pi}{1-\pi},
\qquad
\mathrm{OR}_j=e^{\beta_j}.
\]
Using the notes' threshold \(\tau\) (called \(c\) in the exam),
\[
\hat Y=\bm1\{\hat\pi(\bm X)\geq\tau\}
=\bm1\{\hat p(\bm x)\geq c\}.
\]
The classification metrics are
\[
\begin{aligned}
\mathrm{Accuracy}&=\frac{TP+TN}{TP+TN+FP+FN},&
\mathrm{Precision}&=\frac{TP}{TP+FP},\\
\mathrm{Recall}=\mathrm{TPR}&=\frac{TP}{TP+FN},&
\mathrm{Specificity}&=\frac{TN}{TN+FP}=1-\mathrm{FPR},\\
F_1&=\frac{2(\mathrm{Precision})(\mathrm{Recall})}
{\mathrm{Precision}+\mathrm{Recall}}
=\frac{2TP}{2TP+FP+FN}.&&
\end{aligned}
\]
The cost \(L(c)=4FN(c)+FP(c)\) is problem-specific and is supplied directly
by the exam.

\clearpage
\subsection{Formula-source checklist}

\small
\begin{longtable}{>{\raggedright\arraybackslash}p{0.12\textwidth}
                  >{\raggedright\arraybackslash}p{0.46\textwidth}
                  >{\raggedright\arraybackslash}p{0.31\textwidth}}
\toprule
\rowcolor{navy!10}\textbf{Exam part}&\textbf{Formula used}&\textbf{Course source}\\
\midrule
\endfirsthead
\toprule
\rowcolor{navy!10}\textbf{Exam part}&\textbf{Formula used}&\textbf{Course source}\\
\midrule
\endhead
1(a)&Least-squares criterion, normal equations, and OLS estimator&Chapter 1 notes, p.~2; centered form derived in the solution\\
1(b)&SST--SSR--SSE decomposition, \(R^2\), adjusted \(R^2\), and MSE&Chapter 1 notes, pp.~3--4\\
1(c)&\(\Var(\hat{\bm\beta})=\sigma^2(X^\trans X)^{-1}\)&Chapter 1 notes, p.~4; reparameterization derived in the solution\\
1(d)&Coefficient \(t\) test and partial \(F\) test&Chapter 1 notes, pp.~4--5\\
1(e)&Mean-response CI and prediction interval&Formulas supplied by the exam\\
2(a)&Semi-log linearization and \(e^{\beta}\) interpretation&Chapter 2 notes, p.~5\\
2(b)&Regression MSE, slope SE, and \(t\) interval&Chapter 1 notes, pp.~3--4; quantities supplied by the exam\\
2(c)--(d)&Lognormal mean correction and log-scale prediction interval&Formulas supplied by the exam; back-transformation principle in Chapter 2 notes, pp.~6--7\\
3(a)--(b)&Logit, logistic probability, and odds ratios&Chapter 4 notes, pp.~1--3\\
3(c)&Threshold rule, confusion matrix, and classification metrics&Chapter 4 notes, p.~6\\
3(d)&Weighted empirical cost&Formula supplied by the exam\\
4&OLS covariance and unbiasedness; Gauss--Markov covariance difference&Chapter 1 notes, pp.~2 and 4; proof derived in the solution\\
\bottomrule
\end{longtable}
\normalsize
\end{document}
